\chapter{\texttt{straight.el}: building, locking, and bypassing packages}
\label{ch:straight}
\minitoc\label{straight.el-building-locking-and-bypassing-packages}
\index{straight.el}\index{straight-freeze-versions}\index{lockfile}\index{package management}

% ── Opening signpost ──────────────────────────────────────────────
\texttt{straight.el} replaces \texttt{package.el} with a
clone-build-activate pipeline that makes every package's source
available on disk.  This chapter covers the pipeline mechanics,
the lockfile system, recipe format, autoload caching, and
modification-detection strategies.  Understanding these internals
lets you debug build failures and keep startup fast.
This chapter assumes familiarity with Chapter~\ref{ch:use-package};
the lockfile and recipe concepts introduced here reappear in the
idioms chapter (Chapter~\ref{ch:configuration-idioms}).

\section{The clone-build-activate pipeline}\label{the-clone-build-activate-pipeline}

\texttt{straight.el} replaces \texttt{package.el}'s download-and-activate model with a three-step pipeline applied to every package:

\begin{enumerate}
\item \textbf{Clone}: fetch the source from a Git remote into
  \texttt{\~{}/.emacs.d/\allowbreak straight/\allowbreak repos/\allowbreak <package>/}.
  The lockfile (\texttt{straight/\allowbreak versions/\allowbreak default.el}) records commit
  SHAs, ensuring reproducible installs.
\item \textbf{Build}: byte-compile the source and generate an autoload
  file; outputs go to
  \texttt{\~{}/.emacs.d/\allowbreak straight/\allowbreak build/\allowbreak <package>/}.
  Unlike \texttt{package.el}, straight rebuilds only when the source
  changes (via mtime comparison).
\item \textbf{Activate}: add the build directory to \texttt{load-path}
  and load the generated autoloads.
\end{enumerate}

This means every package is always available as source on disk, so you can edit a package's source directly and \texttt{straight-rebuild-package} to test the change. This is impossible with the \texttt{package.el} model (installed packages are byte-compiled tarballs).

\section{The lockfile and \texttt{straight-freeze-versions}}\label{the-lockfile-and-straight-freeze-versions}

Running \texttt{M-x straight-freeze-versions} writes the current commit SHA of every cloned package to \texttt{$\sim$/.emacs.d/straight/versions/default.el}. Committing this file makes your entire package state reproducible: on a new machine, \texttt{straight-thaw-versions} checks out the exact commits recorded.

The config exempts this file from its \texttt{.gitignore} specifically so it can be versioned:

\begin{minted}{text}
# .gitignore (excerpt)
straight/repos/
straight/build/
# straight/versions/default.el   ← intentionally NOT ignored
\end{minted}

\section{\texttt{:straight nil}: bypassing straight for built-ins}\label{straight-nil-bypassing-straight-for-built-ins}

Not every \texttt{use-package} declaration installs a package from a remote. Built-in packages (\texttt{project}, \texttt{eglot}, \texttt{flymake}, \texttt{treesit}) should not be cloned because they ship with Emacs. \texttt{:straight nil} tells straight to skip the clone-build step and just configure the package:

\begin{minted}{elisp}
(use-package project
  :straight nil          ; built into Emacs, do not clone
  :bind-keymap ("s-p" . project-prefix-map))
\end{minted}

Without \texttt{:straight nil}, straight would attempt to find a recipe for \texttt{project} on MELPA and clone it, potentially shadowing the built-in version.

\section{Git recipe format}\label{git-recipe-format}

For packages not yet on MELPA or requiring a specific fork, straight accepts inline git recipes:

\begin{minted}{elisp}
(use-package treesit-fold
  :straight (treesit-fold :type git
                          :host github
                          :repo "emacs-tree-sitter/treesit-fold"))

(use-package eglot-x
  :straight (eglot-x :type git
                     :host github
                     :repo "nemethf/eglot-x"))
\end{minted}

The recipe overrides whatever MELPA has (or provides a source when MELPA has none). \texttt{:host} can be \texttt{github}, \texttt{gitlab}, or \texttt{codeberg}; \texttt{:repo} is \texttt{<owner>/<name>}. An optional \texttt{:branch} key selects a non-default branch.

\section{\texttt{straight-cache-autoloads} and startup cost}\label{straight-cache-autoloads-and-startup-cost}

By default, straight generates a \textbf{separate} autoload file per package and loads them one by one during \texttt{straight-use-package} calls. With \texttt{straight-cache-autoloads t} (or the equivalent bootstrap option), all autoloads are merged into a single file that is loaded in one \texttt{load} call. The difference can be 30--100 ms on a cold start with many packages.

The bootstrap snippet in this config sets this during the bootstrap phase so it takes effect before any package is added:

\begin{minted}{elisp}
;; In the straight.el bootstrap block (near top of init.el):
(setq straight-cache-autoloads t)
\end{minted}

\section{\texttt{straight-check-for-modifications}: modification detection strategies}\label{straight-check-for-modifications-modification-detection-strategies}

This variable controls how straight decides whether a package needs rebuilding. It must be set \textbf{before the bootstrap snippet}:

\begin{minted}{elisp}
;; Must appear before the bootstrap let-block:
(setq straight-check-for-modifications '(find-when-checking check-on-save))
\end{minted}

The four strategies and their trade-offs are compared in Table~\ref{tab:straight-check-strategies}.

\begin{sidewaystable}[p]
\centering
\caption{Modification check strategies}
\label{tab:straight-check-strategies}
\begin{tabularx}{\linewidth}{@{}llX@{}}
\toprule
Strategy & When it checks & Cost \\
\midrule
\texttt{watch-files} & Continuously via inotify (Linux) / kqueue (macOS) & One FD per repo dir; can exhaust OS limits with 200+ pkgs \\
\texttt{find-when-checking} & On demand when straight explicitly checks a package & Zero at startup; spawns \texttt{find} on check \\
\texttt{check-on-save} & When any \texttt{.el} file is saved (advises \texttt{save-buffer}) & Zero; one lightweight \texttt{advice-add} at startup \\
\texttt{never} & Never & Zero; requires manual \texttt{straight-rebuild-package} \\
\bottomrule
\end{tabularx}
\end{sidewaystable}

The pair \texttt{'(find-when-checking check-on-save)} provides automatic detection from saves plus on-demand scanning, with zero continuous I/O overhead.

\subsection{What happens when no strategy detects a modification}\label{what-happens-when-no-strategy-detects-a-modification}

If straight does not know that a \texttt{.el} file has changed, the stale \texttt{.elc} is loaded silently on the next startup. Force a rebuild explicitly:

\begin{minted}{elisp}
(straight-rebuild-package "package-name")  ;; single package
(straight-rebuild-all)                     ;; all packages (slow)
\end{minted}

% ── Mental model summary ────────────────────────────────────────────
\section{Mental model}
straight.el operates a three-step pipeline for every package: it clones the repository, builds it (byte-compiling and generating autoloads), and then activates it by adding the build directory to \texttt{load-path}.  The lockfile at \texttt{straight/versions/default.el} records exact commit SHAs, so committing it to version control gives fully reproducible installs across machines.

When declaring a built-in package, \texttt{:straight nil} tells straight to skip cloning entirely; without this flag, straight may shadow the bundled version with a stale or different copy.  For forks or packages not registered on MELPA, inline git recipes in the \texttt{use-package} declaration override the default recipe.  Performance-wise, \texttt{straight-cache-autoloads t} merges all package autoloads into a single file, saving 30--100\,ms on cold starts.  The modification-checking strategy \texttt{(find-when-checking check-on-save)} provides automatic rebuild detection from saves and on-demand scanning, with zero continuous I/O overhead.  \(\rightarrow\) See Appendix~\ref{app:straight-internals} for the build-graph resolution algorithm.
